\chapter{Obstacle regularity and quadratic blow-ups}\label{sec:obstacle-blowups}
The classical scalar obstacle problem has two distinct regularity questions: regularity of the minimizing function, and regularity of the boundary of its positivity set. The first is optimally $C^{1,1}$ for a $C^{1,1}$ obstacle. The free boundary can still have singular points. We use \cite{Fig18}, Sections 4--7, and state the deep regularity inputs separately from the subsequent elementary arguments.

\begin{theorem}[Interior optimal regularity for a smooth obstacle problem]\label{thm:obstacle-optimal-regularity}
Let $U\subset\mathbb R^n$, $n\geq2$, be bounded with $C^1$ boundary. Suppose $\psi\in C^{1,1}(\overline U)$ and $g\in C^1(\overline U)$, with $g\geq\psi$ on $\partial U$. The minimizer of $\frac12\int_U|Du|^2$ in the class of Definition \ref{def:obstacle-problem} belongs to $C^{1,1}_{\mathrm{loc}}(U)$.

This is \cite{Fig18}, Theorem 4.1, pp.~5--7, under the smooth-data setup of Section 2. Its proof combines quadratic growth from the contact set with interior harmonic estimates; that full regularity argument remains external here. The conclusion is interior and does not assert $C^2$ regularity or a smooth free boundary. Nonemptiness follows from the competitor $\max\{g,\psi\}$.
\end{theorem}

\begin{lemma}[The Hessian vanishes almost everywhere on contact]\label{lem:contact-hessian-zero}
If $v\in C^{1,1}_{\mathrm{loc}}(U)$ and $v\geq0$, then
\[
 Dv=0\text{ on }\{v=0\},\qquad
 D^2v=0\text{ almost everywhere on }\{v=0\}.
\]
\end{lemma}
\begin{proof}
Every contact point is a local minimum of the differentiable function, so its gradient is zero. The gradient is locally Lipschitz and hence differentiable almost everywhere by Theorem \ref{thm:rademacher}. At almost every contact point $x$, the contact set also has Lebesgue density one by Corollary \ref{cor:lebesgue-density-theorem}. At such a point, for any unit vector $e$ there are contact points $y_r$ with $|y_r-(x+re)|=o(r)$: otherwise a sequence of missing balls of radius comparable to $r$ would contradict density one at $x$. Lipschitz continuity of $Dv$ and $Dv(y_r)=Dv(x)=0$ imply $Dv(x+re)=o(r)$. Differentiability of $Dv$ then gives $D^2v(x)e=0$ in every direction. This proves the claim and supplies the contact-set step in \cite{Fig18}, Section 5, footnote 1.
\end{proof}

\begin{proposition}[The strong equation and complementarity]\label{prop:obstacle-complementarity}
Suppose a minimizer $u$ and its obstacle $\psi$ are $C^{1,1}_{\mathrm{loc}}(U)$. Put $v=u-\psi\geq0$ and $a=-f-\Delta\psi$, with derivatives interpreted almost everywhere. Then
\[
 \Delta v=a\chi_{\{v>0\}},\qquad
 \lambda=a\chi_{\{v=0\}}\,dx\geq0.
\]
In particular $a\geq0$ almost everywhere on contact. This statement assumes the displayed regularity; it does not infer it for arbitrary $L^2$ forcing.
\end{proposition}
\begin{proof}
On $\{v>0\}$ use the noncontact equation. On $\{v=0\}$ the preceding lemma gives $\Delta v=0$ almost everywhere. Because $v$ is $C^{1,1}$ locally, its distributional Laplacian has no singular part. Thus the two identities hold distributionally as well. Finally $\lambda=-\Delta u-f=a-\Delta v$ is nonnegative by Proposition \ref{prop:obstacle-reaction}. Compare \cite{Fig18}, equation (5.1).
\end{proof}

\begin{definition}[The normalized classical obstacle equation]\label{def:normalized-obstacle}
In this and the following section a normalized obstacle solution is a function $u\in C^{1,1}_{\mathrm{loc}}(U)$, $U\subset\mathbb R^n$ open and $n\geq2$, satisfying
\[
 u\geq0,\qquad \Delta u=\chi_{\{u>0\}}\quad\text{almost everywhere in }U.
\]
Its contact set is $\Lambda=\{u=0\}$, its positivity set is $O=\{u>0\}$, and its interior free boundary is $\Gamma=\partial O\cap U$. This is \cite{Fig18}, equation (5.2), after subtracting an obstacle with Laplacian $-1$ in the zero-forcing problem.
\end{definition}

\begin{proposition}[The normalized equation is a variational obstacle equation]\label{prop:normalized-obstacle-minimizes}
On every bounded Lipschitz $B\Subset U$, a normalized obstacle solution minimizes
\[
 \mathcal J(v)=\int_B\left(\tfrac12|Dv|^2+v\right)
\]
among nonnegative $H^1(B)$ functions with $v-u\in H^1_0(B)$. Conversely a local minimizer in this class that is $C^{1,1}_{\mathrm{loc}}$ satisfies the normalized equation.
\end{proposition}
\begin{proof}
For any competitor, integration by parts gives
\[
 \int_B Du\cdot D(v-u)+\int_B(v-u)
 =\int_B(1-\Delta u)(v-u)=\int_{B\cap\Lambda}v\geq0.
\]
The variational inequality therefore proves minimality. Conversely, apply Proposition \ref{prop:obstacle-complementarity} with obstacle zero and forcing $f=-1$. This connects the normalized equation to the preceding variational theory without changing its sign convention.
\end{proof}

\begin{example}[The half-parabola model]\label{ex:obstacle-half-parabola}
For a unit vector $e$, the function $h_e(x)=\frac12((x\cdot e)^+)^2$ is $C^{1,1}$ and solves the normalized equation in all of $\mathbb R^n$. Its gradient is $(x\cdot e)^+e$ and its Hessian is $e\otimes e$ on the positive half-space and zero on the negative half-space. The gradient is continuous across the interface, so the Laplacian has no surface measure. Its free boundary is a hyperplane, and it is generally not $C^2$ there. This is the regular model in \cite{Fig18}, Section 7.1.
\end{example}

\begin{example}[Quadratic models with thin contact sets]\label{ex:obstacle-quadratic-model}
Let $A$ be a symmetric nonnegative $n\times n$ matrix with $\operatorname{tr}A=1$. Then $p_A(x)=\frac12x\cdot Ax$ solves the normalized equation: its contact set is $\ker A$, which has Lebesgue measure zero, and $\Delta p_A=1$ everywhere. Its free boundary is exactly $\ker A$, of dimension between zero and $n-1$. For example, $|x|^2/(2n)$ has an isolated contact point, while $x_1^2/2$ has a hyperplane contact set with positivity on both sides. These differ from the half-parabola despite the last model having the same free-boundary hyperplane. See \cite{Fig18}, Section 7.2.
\end{example}

\begin{proposition}[Quadratic upper growth at the free boundary]\label{prop:obstacle-upper-growth}
Let $u$ solve the normalized equation, let $x_0\in\Gamma$, and suppose $\overline{B_r(x_0)}\subset U$. If $Du$ has Lipschitz constant $L$ on this ball, then
\[
 0\leq\sup_{B_r(x_0)}u\leq\frac L2r^2.
\]
\end{proposition}
\begin{proof}
Continuity and nonnegativity give $u(x_0)=0$ and $Du(x_0)=0$. Along the segment from $x_0$ to $x$, integrate $Du$ and use $|Du(x_0+t(x-x_0))|\leq Lt|x-x_0|$. This gives $u(x)\leq L|x-x_0|^2/2$. This is the Taylor estimate of \cite{Fig18}, Lemma 6.1.
\end{proof}

\begin{proposition}[Quadratic nondegeneracy]\label{prop:obstacle-nondegeneracy}
If $u$ solves the normalized equation, $x_0\in\overline O\cap U$, and $\overline{B_r(x_0)}\subset U$, then
\[
 \sup_{B_r(x_0)}u\geq\frac{r^2}{2n}.
\]
\end{proposition}
\begin{proof}
First suppose $x_0\in O$. The function $h(x)=u(x)-|x-x_0|^2/(2n)$ is harmonic on $O\cap B_r(x_0)$ and is positive at $x_0$. On the portion of its boundary inside $B_r$, it is nonpositive, because that portion is in contact. The maximum principle on the component containing $x_0$ therefore puts a positive maximum on the spherical part of the boundary. There $u>r^2/(2n)$. Continuity equates the supremum over the open ball with the maximum over its closure.

For $x_0\in\Gamma$, take $x_j\in O$ tending to $x_0$ and $r_j=r-|x_j-x_0|>0$. The preceding bound on $B_{r_j}(x_j)\subset B_r(x_0)$ gives $\sup_{B_r(x_0)}u\geq r_j^2/(2n)$. Send $j$ to infinity. This is \cite{Fig18}, Proposition 6.2; the maximum principle can be applied on an arbitrary bounded component because $h$ is continuous on its closure.
\end{proof}

\begin{definition}[Quadratic rescaling and obstacle blow-ups]\label{def:obstacle-rescaling}
For $x_0\in\Gamma$ and $r>0$, set
\[
 u_{x_0,r}(x)=\frac{u(x_0+rx)}{r^2},\qquad x\in(U-x_0)/r.
\]
A blow-up at $x_0$ is a limit of such functions along radii decreasing to zero, in $C^1$ on every compact set. This normalization is dictated by the upper and lower quadratic growth bounds; see \cite{Fig18}, equation (7.1).
\end{definition}

\begin{theorem}[Compactness and nontriviality of obstacle blow-ups]\label{thm:obstacle-blowup-compactness}
Let $u$ be a normalized obstacle solution and $x_0\in\Gamma$. Every sequence $r_j\downarrow0$ has a subsequence for which $u_{x_0,r_j}\to u_0$ in $C^1_{\mathrm{loc}}(\mathbb R^n)$. The limit is a nonzero global normalized obstacle solution, with $0\in\partial\{u_0>0\}$, $u_0(0)=|Du_0(0)|=0$, and a global bound on the Lipschitz constant of its gradient. No homogeneity or classification is assumed in this compactness statement.
\end{theorem}
\begin{proof}
Choose a fixed ball about $x_0$ on which $Du$ is Lipschitz with constant $L$. Scaling gives $D^2u_{x_0,r}(x)=D^2u(x_0+rx)$ almost everywhere, so the same gradient Lipschitz constant works on every fixed ball once $r$ is small. The values and gradients vanish at zero and are consequently bounded on fixed balls. Arzela--Ascoli and a diagonal subsequence give $C^1_{\mathrm{loc}}$ convergence and a $C^{1,1}_{\mathrm{loc}}$ limit with the bound $L$.

Each rescaling satisfies $\Delta u_{x_0,r}=\chi_{\{u_{x_0,r}>0\}}$. Where $u_0>0$, uniform convergence makes the approximants positive locally, so $\Delta u_0=1$ there. On its contact set the Hessian vanishes almost everywhere by Lemma \ref{lem:contact-hessian-zero}. Since the limit has a bounded weak Hessian, these two conclusions give the normalized equation on the whole space. This does not assume convergence of the characteristic functions on the limiting contact set. Nondegeneracy and convergence on every $\overline{B_s}$ give $\sup_{B_s}u_0\geq s^2/(2n)$. Thus the limit is nontrivial and zero remains on its free boundary. This expands the compactness discussion in \cite{Fig18}, pp.~10--11.
\end{proof}

\begin{definition}[Regular free-boundary points]\label{def:obstacle-regular-point}
A point $x_0\in\Gamma$ is regular if some blow-up is $h_e$ from Example \ref{ex:obstacle-half-parabola} for a unit vector $e$. The definition initially requires only one sequence of radii. See \cite{Fig18}, Definition 7.1.
\end{definition}

\begin{definition}[Singular free-boundary points]\label{def:obstacle-singular-point}
A point $x_0\in\Gamma$ is singular if some blow-up is $p_A$ from Example \ref{ex:obstacle-quadratic-model}. Write $\Sigma$ for the set of these points. The definition alone does not establish uniqueness of $A$ or rule out other blow-up types. See \cite{Fig18}, Definition 7.2.
\end{definition}

\begin{theorem}[The classical regular--singular dichotomy]\label{thm:obstacle-dichotomy}
For a normalized classical obstacle solution, every interior free-boundary point belongs to exactly one of the regular and singular classes above. Near each regular point the free boundary is a real-analytic hypersurface consisting only of regular points. Consequently the regular set is relatively open in $\Gamma$ and $\Sigma$ is relatively closed.

This is the classification and regular-point conclusion of \cite{Fig18}, Theorem 7.3, pp.~13--16, attributed there to Caffarelli. Its semiconvexity, classification, and free-boundary improvement arguments remain external. It applies to the constant-coefficient normalized classical obstacle equation, not automatically to thin obstacles, variable forcing, or geometric obstacles.
\end{theorem}
